% ═══════════════════════════════════════════════════════════════════════
\section{Synthèse}
% ═══════════════════════════════════════════════════════════════════════

\begin{frame}{Ce qu'il faut retenir}
  \small
  \begin{itemize}
    \item Un tableau de données est un nuage de points ; on ne le voit
      qu'en projection.
    \item Pythagore : minimiser les distances au sous-espace revient à
      maximiser la variance du nuage projeté.
    \item La variance dans la direction $\vec u$ est $\vec u^{\mathsf T}C\,\vec u$ ;
      ses extremums sont dans les directions propres de la matrice de
      covariance $C$, et valent les valeurs propres.
    \item ACP : diagonaliser $C$ (théorème spectral), ranger
      $\lambda_1\geq\dots\geq\lambda_p$, projeter sur
      $\vec v_1,\dots,\vec v_q$ ; les composantes principales sont non
      corrélées, de variances $\lambda_k$.
    \item La proportion $\frac{\lambda_1+\dots+\lambda_q}{\lambda_1+\dots+\lambda_p}$
      dit jusqu'à quel point le dessin est fidèle ; l'erreur de
      reconstruction moyenne est $\lambda_{q+1}+\dots+\lambda_p$.
    \item Réduire (matrice de corrélation) dès que les unités ou les
      ordres de grandeur diffèrent ; se méfier des valeurs aberrantes.
    \item Interpréter les axes à l'aide des composantes des vecteurs
      propres : « niveau général », « taille du bec », « performance ».
  \end{itemize}
\end{frame}

\begin{frame}{Des notions du secondaire aux applications modernes}
  \small
  \begin{center}
  \begin{tabular}{ll}
    \toprule
    Notion du secondaire & Rôle dans l'ACP \\
    \midrule
    nuage de points, corrélation & le point de départ ; la forme du nuage \\
    droite de régression & un point de comparaison avec l'axe principal \\
    cote $z$ & réduire les variables \\
    fonction sinusoïdale & la variance selon l'angle d'une droite \\
    vecteurs, produit scalaire & projeter \\
    théorème de Pythagore & l'équivalence des deux problèmes \\
    \bottomrule
  \end{tabular}
  \end{center}
  \vspace{6pt}
  Une méthode de 1901, devenue l'un des outils les plus employés de la
  statistique et de la science des données : le calcul de quelques vecteurs
  propres d'une matrice symétrique.
\end{frame}

\begin{frame}{Et maintenant}
  \begin{itemize}
    \item Les démonstrations laissées « au tableau » sont toutes dans les
      notes de cours, avec tous les détails.
    \item Exercices du chapitre (avec solutions complètes) : calculs à la
      main (5.1 à 5.6), dont les valeurs extrêmes (5.5) ; programmes Python
      (5.7 à 5.11), dont l'erreur de saisie (5.10) ; réflexion pour la
      classe (5.12).
    \item L'application interactive des manchots est sur le site des notes :
      \href{\urlappli}{\color{sectioncolor}ouvrir l'application}.
    \item Les rappels d'algèbre linéaire (projection orthogonale, valeurs
      propres, théorème spectral) sont à l'annexe D.
  \end{itemize}
\end{frame}
